\section{Discussion}
\label{sec:discussion}

TIGER profiles provide a vocabulary for design conversations that quality scores cannot.
The finding that $\geq 90\%$ of expert-rated game pairs are correctly discriminated by
TIGER suggests that the five dimensions capture distinctions that experienced players
already make intuitively but struggle to articulate. TIGER externalizes that intuition
into a formal, reproducible scoring scheme.

\subsection{TIGER vs. Quality Frameworks}

TIGER does not replace quality assessment: BGG ratings and similar quality signals
remain the primary tool for evaluating whether a game is worth playing. TIGER supplements
quality by characterizing \emph{what kind} of game it is. The two frameworks answer
different questions and should be used together, not substituted for one another.

\subsection{Limitations}

The anchor descriptors at 0, 5, and 10 are derived from the TIGRIS corpus, which is
biased toward Euro-tradition games published 1995--2024. The Range dimension in particular
may not generalize well to genres not well-represented in the corpus (dexterity games,
legacy games, campaign games). The expert panel validation used five evaluators, which
provides useful signal but is underpowered for detecting differential accuracy across
dimensions. A larger human calibration study is future work.

\subsection{Practical Applications}

TIGER profiles are immediately useful for game recommendation (find games with similar
profiles to a known favorite), design critique (identify which dimensions a prototype
under-develops relative to its genre), and market analysis (identify which TIGER regions
are over-crowded versus under-served).
